% chuong_v3/fig_ch08_window_effect.tex
% Minh họa hiệu ứng suy giảm hình học của hàm tự tương quan cửa sổ Hanning Rw[tau]
% và cơ chế bù suy hao qua hệ số C[tau] = round(2^14 * Rw[0] / Rw[tau])
\begin{tikzpicture}[x=0.06cm, y=4.5cm, >=Stealth]
    % Lưới nền mờ
    \draw[very thin, gray!20, step=1] (0,-0.5) grid[xstep=30, ystep=0.2] (225,1.2);

    % Trục tọa độ
    \draw[->, thick, black!70] (0,0) -- (235,0) node[right, font=\footnotesize] {Độ trễ $\tau$ (mẫu)};
    \draw[->, thick, black!70] (0,-0.55) -- (0,1.25) node[above, font=\footnotesize] {Biên độ chuẩn hóa};

    % Vạch chia trục hoành
    \foreach \x in {0, 26, 60, 100, 140, 180, 214} {
        \draw[thick, black!60] (\x,0.02) -- (\x,-0.02) node[below, font=\scriptsize, fill=white, inner sep=1pt] {\x};
    }
    % Vạch chia trục tung
    \foreach \y in {-0.4, -0.2, 0.0, 0.2, 0.4, 0.6, 0.8, 1.0} {
        \draw[thick, black!60] (2,\y) -- (-2,\y) node[left, font=\scriptsize] {\y};
    }

    % Vùng quét trễ tau in [26, 214]
    \fill[blue!5, opacity=0.7] (26,-0.52) rectangle (214,1.16);
    \node[blue!70!black, font=\scriptsize, fill=white, inner sep=2.5pt, rounded corners=2pt, draw=blue!30] at (120,1.16) 
        {Dải quét trễ phần cứng: $\tau \in [26, 214]$ ($F_0 \in [75, 600]\text{ Hz}$)};

    % Đường cong 1: Tự tương quan của cửa sổ Hanning Rw[tau]/Rw[0]
    \draw[thick, red!80!black, smooth] plot coordinates {
        (0,1.000) (15,0.996) (26,0.989) (40,0.973) (60,0.943) (80,0.902) 
        (100,0.851) (120,0.791) (140,0.725) (160,0.654) (180,0.584) (200,0.514) (214,0.468)
    };
    % Chú thích Rw[tau] đặt ở vùng thoáng (tau in [25, 65], y in [0.45, 0.75]) với mũi tên chỉ lên đường cong đỏ
    \node[red!80!black, font=\scriptsize, fill=white, inner sep=2.5pt, rounded corners=1pt, draw=red!40] (lbl_rw) at (46,0.60) 
        {$R_w[\tau] / R_w[0]$ (Cửa sổ co ngắn dần)};
    \draw[->, thick, red!80!black] (lbl_rw.north) -- (46,0.965);

    % Đường cong 2: Tự tương quan chưa chuẩn hóa Rxw[tau]/Rxw[0]
    \draw[thick, gray!60!black, dashed, smooth] plot coordinates {
        (0,1.000) (20,0.05) (40,-0.45) (60,-0.30) (80,0.25) (95,0.80) 
        (100,0.851) (105,0.80) (120,0.15) (140,-0.35) (160,-0.20) (180,0.15) 
        (195,0.48) (200,0.514) (205,0.48) (214,0.20)
    };
    % Chú thích Rxw đặt ở khoảng thoáng trung tâm (tau in [122, 154], y in [0.05, 0.19]), mũi tên chỉ thẳng xuống đáy đường cong xám
    \node[gray!70!black, font=\scriptsize, align=center, fill=white, inner sep=2.5pt, rounded corners=1pt, draw=gray!40] (lbl_rxw) at (138,0.12) 
        {$R_{xw}[\tau] / R_{xw}[0]$\\(Chưa chuẩn hóa: đỉnh bị dìm)};
    \draw[->, thick, gray!60!black] (lbl_rxw.south) -- (138,-0.33);

    % Đường cong 3: Sau khi chuẩn hóa U[tau] / R0n
    \draw[very thick, blue!85!black, smooth] plot coordinates {
        (0,1.000) (20,0.05) (40,-0.46) (60,-0.32) (80,0.28) (95,0.94) 
        (100,1.000) (105,0.94) (120,0.19) (140,-0.48) (160,-0.31) (180,0.26) 
        (195,0.93) (200,1.000) (205,0.93) (214,0.43)
    };
    \node[blue!85!black, font=\scriptsize, fill=white, inner sep=2pt, rounded corners=1pt] at (75,1.06) 
        {Đỉnh $\tau_0 = 100$ phục hồi $1{,}0$};
    \node[blue!85!black, font=\scriptsize, fill=white, inner sep=2pt, rounded corners=1pt] at (175,1.06) 
        {Đỉnh $2\tau_0 = 200$ phục hồi $1{,}0$};

    % Điểm mốc và mũi tên chỉ thị
    \filldraw[red!80!black] (214,0.468) circle (1.5pt);
    \node[above right, red!80!black, font=\scriptsize, fill=white, inner sep=1.5pt] at (214,0.50) 
        {$0{,}468$ tại $\tau=214$};

    \filldraw[blue!85!black] (100,1.0) circle (1.5pt);
    \filldraw[blue!85!black] (200,1.0) circle (1.5pt);
    \draw[->, thick, green!60!black] (100,0.851) -- (100,0.98);
    \draw[->, thick, green!60!black] (200,0.514) -- (200,0.98);

    \node[right, green!50!black, font=\scriptsize, fill=white, inner sep=1.5pt] at (102,0.92) {Bù $\times 1{,}175$};
    \node[left, green!50!black, font=\scriptsize, fill=white, inner sep=1.5pt] at (197,0.75) {Bù $\times 1{,}946$};
\end{tikzpicture}
